\scanpage{13}

Pour \deleted{que} une $\deleted{\unsure{\varphi|C}}C$-polarisation donnée, les formes $\varphi'$ sur le $=$ $M$ qui sont des $C$-polarisations, sont \uline{exactement} celles qui sont compatibles à $\varphi$.

c) Donc : Supposons $\varepsilon(C)=1$. Soit, pour tt $M$ donné, $P(M)$ l'ens de ses $C$-polarisations. Alors $M\mapsto P(M)$ est une polarisation de $\mathcal M$, \deleted{\unsure{à|et}} cela près que les $P(M)$ peuvent ne pas être tous $\ne\varnothing$.

16) \uline{Questions} Toute polarisation de \deleted{\unsure{$\mathcal M$}} $\mathcal M=\operatorname{Rep}(G)$ peut-elle se décrire par un $C$ comme dessus ($C^2=\unsure{y|\gamma|\eta},\ \varepsilon(C)=1$) ? Est-ce vrai si $G$ connexe ? Cas $G'$ \inserted{$=\operatorname{Ker}\varepsilon$} compact ?

\margin{$?$}

\uline{NB} Un tel $C$ \inserted{mod centre $Z$} est uniquement déterminé mod \deleted{la} \deleted{\unsure{centralité}} mod conjugaison, \deleted{\unsure{pour éléments}} du moins si $G$ est connexe. Donc \deleted{\unsure{l'on|bien}} l'indétermination de $C$ \inserted{dans ce cas} est modulo points d'ordre $2$ de $Z(\mathbb R)$ et conjugaison.

\uline{Re NB} Pour qu'il existe un $C$ comme dessus fournissant une polarisation, il \inserted{cf. 14} suffit, en termes du groupe $G'=\operatorname{Ker}\varepsilon$ muni de $y=i(-1)$, qu'il existe un $C\in G'(\mathbb R)$ tel que \uline{\unsure{toute|toutes}} \unsure{repr|reprs|rep.} \inserted{de $G'$} sur laquelle $y$ induit $\pm1$ admette une $C$-polarisation, i.e. une forme \inserted{inv.} $\varphi$ telle que $\varphi(x,Cy)$ soit sym. déf $>0$.

On dit que le groupe \inserted{alg affine} réel $G'$ est \uline{hodgien} s'il existe un tel $C$ \deleted{d'ordre 4} \inserted{avec $C^4=1$} ($y=C^2$ n'étant pas fixé à priori).

% archive sheet number: 12
